\subsection{The actual tensor and conjugate-extension comparisons}

\begin{result}{A.4.36}{Symmetric-square tensor and conjugate-extension map}{Proved}
For a real algebraic short model with three real roots, choose
$\omega_1=w>0$ and $\omega_2=\tau w\in i\mathbf R_{>0}$.
Put
\[
 T_a=\begin{pmatrix}2a&1\\1&0\end{pmatrix},\qquad
 P=\begin{pmatrix}w&-(h+c)w\\ \tau w&-(h-c)\tau w\end{pmatrix}.
\]
The actual K\"unneth class
\[
 \theta_a=\operatorname{pr}_1^*\omega\wedge\operatorname{pr}_2^*\xi
 +\operatorname{pr}_1^*\xi\wedge\operatorname{pr}_2^*\omega
 +2a\operatorname{pr}_1^*\omega\wedge\operatorname{pr}_2^*\omega
\]
has product-cycle period matrix $PT_aP^{\mathsf T}$, and
\[
 \frac{PT_aP^{\mathsf T}}{2\pi i}
 =\begin{pmatrix}
 (a-h-c)/(c\tau)&(a-h)/c\\
 (a-h)/c&(a-h+c)\tau/c
 \end{pmatrix}.
\]
At $a=h$ this matrix is $\operatorname{diag}(-\tau^{-1},\tau)$.
For a non-CM curve no nonzero complex scalar multiple of this Betti
vector has all coordinates algebraic. The determinant-twisted
symmetric-square comparison identifies this class with the earlier
Hodge projector by an explicit isomorphism.

For every algebraic short model, without the reality restriction,
the unique complex-linear map carrying its universal-extension
period basis to the conjugate period basis is
\[
 M=\begin{pmatrix}
 -h/c&1/c\\ c-h\bar h/c&\bar h/c
 \end{pmatrix}.
\]
It descends to a holomorphic group isomorphism
$(E^\dagger)^{\mathrm{an}}\to(\bar E^\dagger)^{\mathrm{an}}$, with
$\det M=-1$ and $\bar M M=I$. This map is not algebraic; its entries
cannot all be algebraic.
\end{result}

\begin{proof}
A.4.1 gives
$\eta_1=(h+c)w$, $\eta_2=(h-c)\tau w$ and
$c\tau w^2=\pi i$. Thus $\det P=2\pi i$, and multiplication gives
the displayed tensor matrix. The class $\theta_h$ has type $(1,1)$,
as $\xi+h\omega$ is antiholomorphic. Its two nonzero normalized
Betti coordinates have ratio $-\tau^2$. Schneider's theorem,
in the form proved in A.3.1, makes $\tau$ transcendental for this
non-CM algebraic curve; algebraicity of $\tau^2$ would make $\tau$
algebraic. Algebraicity of both coordinates after any nonzero
rescaling is therefore impossible.

To specify the projector comparison, put
$J=\left(\begin{smallmatrix}0&1\\-1&0\end{smallmatrix}\right)$.
The map
\[
 \operatorname{Sym}^2 V\otimes(\det V)^{-1}
       \longrightarrow\mathfrak{sl}(V),\qquad T\longmapsto TJ
\]
is an isomorphism: its image has trace zero, and $-UJ$ is symmetric
for every trace-zero $U$. The identity
$P^{\mathsf T}JP=(\det P)J$ gives its exact compatibility with
comparison:
\[
 (PTP^{\mathsf T}/\det P)J=P(TJ)P^{-1}.
\]
In particular
\[
 T_hJ=\begin{pmatrix}-1&2h\\0&1\end{pmatrix}=2e_h-I,
 \qquad e_h=\begin{pmatrix}0&h\\0&1\end{pmatrix},
\]
and the Betti matrix is
$\left(\begin{smallmatrix}0&-\tau^{-1}\\-\tau&0\end{smallmatrix}\right)$.
This is the same projector as A.4.18 through the displayed morphism.
Its de Rham centralizer is exactly
\[
 \left\{\begin{pmatrix}r&a(s-r)\\0&s\end{pmatrix}:
                  r,s\in\mathbf C^\times\right\}
 \quad\text{for }T_aJ.
\]
Indeed a commuting matrix preserves its two distinct eigenlines
$\mathbf C\omega$ and $\mathbf C(\xi+a\omega)$ and acts there by
$r,s$. Consequently an algebraic two-dimensional stabilizer exists
for every algebraic complement $a$; the Betti stabilizer is obtained
by the actual conjugation $C_a\mapsto PC_aP^{-1}$.
The comparison coordinates just computed fail the algebraic-Betti
hypothesis in \cite[Definition 4.5, p.13; Theorem 7.12(ii), p.27 of
 the opened snapshot]{EllKSVAuthor45} for this proposed tensor.

For the general conjugate-extension assertion, write a vector as
$(u,hu+cv)$. The displayed $M$ sends it to
$(v,\bar h v+cu)$. Taking $u=\omega_i$, $v=\bar\omega_i$ proves
that it sends the two actual periods to their conjugates.
Those periods are a complex basis by Legendre, proving uniqueness.
Direct substitution proves $\det M=-1$ and $\bar M M=I$.
In the shifted coordinates $(u,s-hu)$ and $(v,t-\bar h v)$, the same
map has matrix
$\left(\begin{smallmatrix}0&1/c\\c&0\end{smallmatrix}\right)$.

Every algebraic group map $E^\dagger\to\bar E^\dagger$ sends the
vertical $\mathbf G_a$ into the vertical $\mathbf G_a$.
Its composition with $\bar E^\dagger\to\bar E$ is constant:
a morphism $\mathbf A^1\to\bar E$ extends to $\mathbf P^1$ by
properness and is constant by Riemann--Hurwitz. But $M(0,1)$ has
first coordinate $1/c\ne0$, proving nonalgebraicity.
If every entry of $M$ were algebraic, the period-preserving Lie-map
theorem \cite[Theorem 3.1, p.13 of the opened preprint]{EllBostCharles2016}
would make this map algebraic: the source period lattice spans its
complex Lie algebra. This gives the final assertion.
\end{proof}

The matrix computations are reproduced by
\texttt{check\_symmetric\_hodge\_tensor.py}. The remaining scale $c$
is present in the explicit comparison. These results supply no
non-CM exclusion for algebraic $h$ alone.
